% Options for packages loaded elsewhere
\PassOptionsToPackage{unicode}{hyperref}
\PassOptionsToPackage{hyphens}{url}
\documentclass[
]{article}
\usepackage{xcolor}
\usepackage[margin=1in]{geometry}
\usepackage{amsmath,amssymb}
\setcounter{secnumdepth}{-\maxdimen} % remove section numbering
\usepackage{iftex}
\ifPDFTeX
  \usepackage[T1]{fontenc}
  \usepackage[utf8]{inputenc}
  \usepackage{textcomp} % provide euro and other symbols
\else % if luatex or xetex
  \usepackage{unicode-math} % this also loads fontspec
  \defaultfontfeatures{Scale=MatchLowercase}
  \defaultfontfeatures[\rmfamily]{Ligatures=TeX,Scale=1}
\fi
\usepackage{lmodern}
\ifPDFTeX\else
  % xetex/luatex font selection
\fi
% Use upquote if available, for straight quotes in verbatim environments
\IfFileExists{upquote.sty}{\usepackage{upquote}}{}
\IfFileExists{microtype.sty}{% use microtype if available
  \usepackage[]{microtype}
  \UseMicrotypeSet[protrusion]{basicmath} % disable protrusion for tt fonts
}{}
\makeatletter
\@ifundefined{KOMAClassName}{% if non-KOMA class
  \IfFileExists{parskip.sty}{%
    \usepackage{parskip}
  }{% else
    \setlength{\parindent}{0pt}
    \setlength{\parskip}{6pt plus 2pt minus 1pt}}
}{% if KOMA class
  \KOMAoptions{parskip=half}}
\makeatother
\usepackage{graphicx}
\makeatletter
\newsavebox\pandoc@box
\newcommand*\pandocbounded[1]{% scales image to fit in text height/width
  \sbox\pandoc@box{#1}%
  \Gscale@div\@tempa{\textheight}{\dimexpr\ht\pandoc@box+\dp\pandoc@box\relax}%
  \Gscale@div\@tempb{\linewidth}{\wd\pandoc@box}%
  \ifdim\@tempb\p@<\@tempa\p@\let\@tempa\@tempb\fi% select the smaller of both
  \ifdim\@tempa\p@<\p@\scalebox{\@tempa}{\usebox\pandoc@box}%
  \else\usebox{\pandoc@box}%
  \fi%
}
% Set default figure placement to htbp
\def\fps@figure{htbp}
\makeatother
\setlength{\emergencystretch}{3em} % prevent overfull lines
\providecommand{\tightlist}{%
  \setlength{\itemsep}{0pt}\setlength{\parskip}{0pt}}
\usepackage{bookmark}
\IfFileExists{xurl.sty}{\usepackage{xurl}}{} % add URL line breaks if available
\urlstyle{same}
\hypersetup{
  pdftitle={laporan},
  hidelinks,
  pdfcreator={LaTeX via pandoc}}

\title{laporan}
\author{}
\date{\vspace{-2.5em}}

\begin{document}
\maketitle

\section{BAB I PENDAHULUAN}\label{bab-i-pendahuluan}

\subsection{1.1 Latar Belakang}\label{latar-belakang}

Pembangunan wilayah yang komprehensif tidak hanya diukur dari
pertumbuhan infrastruktur fisik, tetapi juga dari peningkatkan taraf
hidup dan kesejahteraan masyarakatnya. Salah satu indikator utama yang
masih menjadi tantangan dalam pembangunan di berbagai wilayah di
Indonesia adalah tingkat kemiskinan (Kuncoro, 2010). Kemiskinan
merupakan fenomena multidimensional yang dipengaruhi oleh berbagai
faktor sosial dan ekonomi. Dalam pendekatan makroekonomi, kualitas
sumber daya manusia dan ketersediaan lapangan pekerjaan dianggap sebagai
dua pilar krusial yang dapat mengentaskan masyarakat dari garis
kemiskinan.

Teori human capital atau modal manusia menegaskan bahwa pendidikan
formal merupakan investasi jangka panjang yang berbanding lurus dengan
peningkatan produktivitas dan kapabilitas ekonomi seseorang (Todaro \&
Smith, 2015). Hal ini direpresentasikan oleh indikator rata-rata lama
sekolah. Logika statistiknya, semakin tinggi durasi pendidikan yang
diselesaikan oleh penduduk di suatu wilayah, semakin besar peluang
mereka untuk diserap oleh pasar kerja yang layak, sehingga tingkat
kemiskinan dapat ditekan. Sebaliknya, ketidakmampuan pasar kerja dalam
menyerap angkatan kerja akan memicu tingginya angka pengangguran
terbuka. Tingkat pengangguran yang tinggi secara langsung memutus aliran
pendapatan individu, yang pada akhirnya bermuara pada peningkatan
persentase kemiskinan wilayah (Sukirno, 2011).

Dalam pengolahan data riil terkait indikator-indikator tersebut, analis
sering dihadapkan pada masalah kualitas data, seperti adanya data yang
hilang (missing value) atau nilai anomali yang ekstrem (outlier). Oleh
karena itu, proyek ini disusun untuk melakukan analisis statistika
secara komprehensif menggunakan bahasa pemrograman R. Evaluasi
difokuskan pada penerapan statistika deskriptif, penanganan
ketidaksempurnaan data, hingga pembuktian statistik mengenai korelasi
antara rata-rata lama sekolah dan tingkat pengangguran terhadap
persentase kemiskinan di berbagai wilayah.

\subsection{1.2 Tujuan}\label{tujuan}

Melalui proyek akhir ini mahasiswa diharapkan mampu:

\begin{enumerate}
\def\labelenumi{\arabic{enumi}.}
\tightlist
\item
  Memahami struktur dan karakteristik dataset.
\item
  Melakukan eksplorasi data menggunakan R.
\item
  Mengidentifikasi missing value dan outlier.
\item
  Melakukan statistika deskriptif.
\item
  Membuat visualisasi data yang informatif.
\item
  Menggunakan konsep probabilitas dan distribusi data.
\item
  Melakukan analisis hubungan antar variabel.
\item
  Menarik kesimpulan berbasis data.
\item
  Menyusun laporan ilmiah sederhana.
\item
  Mengkomunikasikan hasil analisis secara sistematis.
\end{enumerate}

\subsection{1.3 Deskripsi Dataset}\label{deskripsi-dataset}

Dataset yang digunakan dalam analisis ini adalah
pembangunan\_wilayah\_missing\_outlier.csv. Untuk menjawab tujuan
pengamatan, analisis difokuskan pada dua variabel utama pengukur
kualitas sumber daya manusia dan tenaga kerja, beserta variabel
identitas wilayahnya. Berikut adalah rincian variabel yang digunakan:

\begin{enumerate}
\def\labelenumi{\arabic{enumi}.}
\tightlist
\item
  wilayah (Character):
\item
  provinsi (Character):
\item
  tahun (Integer):
\item
  pdrb\_perkapita (Numeric):
\item
  kemiskinan (Numeric):
\item
  pengangguran (Numeric):
\item
  ipm (Numeric):
\item
  harapan\_hidup (Numeric):
\item
  rata\_lama\_sekolah (Numeric):
\item
  akses\_internet (Numeric):
\item
  jalan\_baik (Numeric):
\item
  air\_bersih (Numeric):
\item
  catatan\_data (Character):
\end{enumerate}

\section{BAB II METODOLOGI}\label{bab-ii-metodologi}

\subsection{2.1 Tools}\label{tools}

\begin{itemize}
\tightlist
\item
  R dan R Studio: Bahasa pemrograman dan lingkungan IDE (Integrated
  Development Environment) yang digunakan untuk manipulasi dataframe,
  ekstraksi nilai statistik deskriptif, hingga pembuatan plot dan
  visualisasi hubungan antar variabel.
\item
  GitHub: Platform version control yang digunakan sebagai wadah
  kolaborasi kelompok untuk menyimpan, melacak perubahan, dan menyatukan
  script laporan analisis secara terpusat.
\end{itemize}

\subsection{2.2 Metode Analisis}\label{metode-analisis}

Analisis laporan ini berjalan melalui serangkaian proses komputasi
statistik dengan tahapan berikut:

\section{BAB III HASIL DAN
PEMBAHASAN}\label{bab-iii-hasil-dan-pembahasan}

\begin{center}\rule{0.5\linewidth}{0.5pt}\end{center}

\subsection{3.1 Statistika Deskriptif}\label{statistika-deskriptif}

\subsection{\texorpdfstring{3.2 \emph{Missing
Value}}{3.2 Missing Value}}\label{missing-value}

Berdasarkan pemeriksaan awal, ditemukan adanya data yang hilang pada
beberapa variabel. Untuk mengatasi hal ini tanpa mengurangi ukuran
sampel, dilakukan \emph{median imputation} atau mengisi nilai yang
kosong menggunakan nilai tengah dari masing-masing variabel. Hasil akhir
menunjukkan seluruh kolom telah bersih dari nilai \emph{missing value}.

\subsection{\texorpdfstring{3.3
\emph{Outlier}}{3.3 Outlier}}\label{outlier}

Berdasarkan visualisasi menggunakan boxplot pada variabel pengangguran
dan rata-rata lama sekolah, terdeteksi adanya beberapa outlier.
Pembersihan dilakukan dengan menggunakan metode IQR (\emph{Interquartile
Range}) dengan ambang batas 1.5 kali nilai IQR. Dataset lalu diperbarui,
dimana data di luar rentang batas bawah dan batas atas sudah disaring.

\subsection{3.4 Visualisasi}\label{visualisasi}

\subsection{3.5 Distribusi Data}\label{distribusi-data}

\subsection{3.6 Korelasi}\label{korelasi}

\subsection{3.7 Interpretasi}\label{interpretasi}

\section{BAB IV KESIMPULAN}\label{bab-iv-kesimpulan}

\begin{center}\rule{0.5\linewidth}{0.5pt}\end{center}

\section{LAMPIRAN}\label{lampiran}

\begin{center}\rule{0.5\linewidth}{0.5pt}\end{center}

\end{document}
